
\section{Setting the stage}
We review several definitions and results on simplicial complexes, Stanley--Reisner rings, stress spaces, and Cohen--Macaulayness, as well as prepare ground for the proofs. For all undefined terminology we refer the reader to~\cite{Lee96, Stanley96}.

A(n abstract) \emph{simplicial complex} $\Delta$ on the ground set $V$ is a collection of subsets of $V$ that is closed under inclusion; $v$ is a \emph{vertex} of $\Delta$ if $\{v\}\in \Delta$, but not all elements of $V$ are required to be vertices. The elements of $\Delta$ are called \emph{faces}. The \emph{dimension of a face} $\tau\in\Delta$ is $\dim\tau\coloneqq |\tau|-1$. The \emph{dimension of $\Delta$}, $\dim\Delta$, is the maximum dimension of its faces. A face of a simplicial complex $\Delta$ is a \emph{facet} if it is maximal w.r.t.~inclusion. We say that $\Delta$ is \emph{pure} if all facets of $\Delta$ have the same dimension. To simplify notation, for a face that is a vertex, we write $v$ instead of $\{v\}$; we also define the following two subcomplexes of $\Delta$ called the \emph{star of $v$} and the \emph{link of $v$ in $\Delta$}:
$\st_\Delta(v)=\st(v):=\{\sigma\in\Delta \ : \ \sigma\cup v\in\Delta\}$ and $\lk_\Delta(v)=\lk(v)\coloneqq  \{\sigma\in \st_\Delta(v) \ : \ v\notin\sigma\}$.


Let $\Delta$ be a $(d-1)$-dimensional simplicial complex. For $-1 \leq i \leq d-1$, the $i$th \emph{$f$-number} of $\Delta$, $f_i = f_i(\Delta)$, denotes the number of $i$-dimensional faces of $\Delta$. The \emph{$h$-numbers} of $\Delta$, $h_i=h_i(\Delta)$ for $0 \leq i \leq d$, are defined by the relation $\sum_{i=0}^{d}h_i\lambda^{d-i}=\sum_{i=0}^{d}f_{i-1}(\lambda-1)^{d-i}$. Finally, the \emph{$g$-numbers} of $\Delta$ are $g_0(\Delta)\coloneqq 1$ and $g_i(\Delta)\coloneqq h_i(\Delta)-h_{i-1}(\Delta)$ for $1 \leq i \leq d/2$.



Let $\Delta$ be a simplicial complex on the ground set $V$. Let $X = \{x_v : v \in V\}$ be the set of variables and let $\R[X]$ be the polynomial ring over the real numbers $\R$ in variables $X$. The \emph{Stanley--Reisner ideal} of $\Delta$ is defined as 
\[
I_\Delta=\left(x_{v_1}x_{v_2}\dots x_{v_i}: \ \{v_1, v_2, \dots, v_i\}\notin \Delta\right),
\] 
\ie
it is the ideal generated by the squarefree monomials corresponding to non-faces of $\Delta$. The \emph{Stanley--Reisner ring} of $\Delta$ is $\R[\Delta] \coloneqq  \R[X]/I_\Delta$. The ring $\R[\Delta]$ has an $\N$-grading: $\R[\Delta]= \bigoplus_{i=0}^{\infty} \R[\Delta]_i$, where the $i$th graded component $\R[\Delta]_i$ is the space of homogeneous elements of degree $i$ in $\R[\Delta]$. In general, for an $\N$-graded vector space $M$, denote by $M_i$ the $i$th graded component of~$M$.

Let $\Delta$ be a simplicial complex and let $\Theta=\theta_1,\ldots,\theta_\ell$ be a sequence of linear forms in $\R[X]$, where $\ell$ is a nonnegative integer. Denote the quotient $\R[\Delta]/\Theta\R[\Delta]$ by $\R(\Delta,\Theta)$. 

For our proofs, we will work in the dual setting of stress spaces developed by Lee~\cite{Lee96}, see also~\cite[Section 3]{Adiprasito-g-conj}. It should also be mentioned that stress spaces are essentially the same objects as \emph{inverse systems} in commutative algebra --- the notion that goes back to Macaulay; see~\cite[Theorem 21.6 and Exercise 21.7]{Eisenbud}. Observe that a variable $x_v$ acts on $\R[X]$ by $\frac{\partial}{\partial{x_v}}$; for brevity, we will denote this operator by $\partial_{x_v}$. More generally, if $c(X)=\sum_{v\in V} c_v x_v$ is a linear form in $\R[X]$, then we define \begin{equation*}
\begin{split}
\partial_{c(X)} : \R[X]&\to\R[X], \\
w&\mapsto \sum_{v\in V}c_v\cdot\partial_{x_v}w=\sum_{v\in V}c_v\frac{\partial w}{\partial{x_v}}.
\end{split}
\end{equation*}
	
For a monomial $\mu\in \R[X]$, the \emph{support} of $\mu$ is $\supp(\mu)=\{v\in V : x_v\,|\,\mu\}$. A homogeneous polynomial $w\in\R[X]$ of degree $i$ is called an \emph{$i$-stress} on $\Delta$ w.r.t.~$\Theta=\theta_1,\ldots,\theta_\ell$ if it satisfies the following conditions:
	\begin{itemize}
	\item Every term $\mu$ of $w$ is supported on a face of $\Delta$: $\supp(\mu)\in\Delta$, and
	\item $\partial_{\theta_k}w=0$ for all $k=1,\ldots,\ell$.
	\end{itemize}
The \emph{support} of an $i$-stress $w$, $\supp (w)$, is the subcomplex of $\Delta$ generated by the support of all terms of $w$. We say that a face $F\in\Delta$ \emph{participates in a stress $w$} if $F\in \supp(w)$. We also say that a stress $w$ \emph{lives on a subcomplex $\Gamma$} of $\Delta$ if $\supp(w)\subseteq \Gamma$.

Denote the set of all $i$-stresses on $\Delta$ w.r.t.~$\Theta$ by $\Stress(\Delta, \Theta)_i$. 
This set is a vector space~\cite{Adiprasito-g-conj,Lee96}; it is a subspace of $\R[X]$. In fact, $\Stress(\Delta, \Theta)_i$ is the orthogonal complement of $(I_\Delta+(\Theta))_i$ in $\R[X]_i$ w.r.t.~a certain inner product on $\R[X]_i$, see~\cite[Section 3]{Lee96}. Thus, as a vector space, $\Stress(\Delta, \Theta)_i$ is canonically isomorphic to $\R(\Delta, \Theta)_i$. (For an alternative approach using the Weil duality, see~\cite[Section 3]{Adiprasito-g-conj}.) Another very useful and easy fact is that for every linear form $c(X)\in \R[X]$, the operator $\partial_{c(X)}$ maps $\Stress(\Delta,\Theta)_i$ into $\Stress(\Delta, \Theta)_{i-1}$, that is, if $w$ is a stress, then so is $\partial_{c(X)}w$. This follows from the fact that $\partial_{\theta_k}$ and $\partial_{c(X)}$ commute, and that a subset of a face of $\Delta$ is a face of $\Delta$.

Stresses are convenient to work with for the following reason: if $\Gamma$ is a subcomplex of $\Delta$ (considered as a complex on the same ground set $V$ as $\Delta$), then there is a natural surjective homomorphism $\rho:\R[\Delta] \to \R[\Gamma]$; it induces a surjective homomorphism $\R(\Delta,\Theta) \to \R(\Gamma,\Theta)$. On the level of stress spaces, the situation is much easier to describe: $\Stress(\Gamma,\Theta)_i$ is a subspace of $\Stress(\Delta,\Theta)_i$.

A simplicial complex $\Delta$ is \emph{centrally symmetric} or \emph{cs} if its ground set is endowed with a \emph{free involution} $\alpha: V \rightarrow V$ that induces a free involution on the set of all non-empty faces of $\Delta$. In more detail, for all non-empty faces $\tau \in \Delta$, the following holds: $\alpha(\tau)\in\Delta$, $\alpha(\tau)\neq \tau$, and $\alpha(\alpha(\tau))=\tau$. To simplify notation, we write $\alpha(\tau)=-\tau$ and refer to $\tau$ and $-\tau$ as \emph{antipodal faces} of $\Delta$. 

A large family of cs simplicial complexes is given by cs simplicial polytopes. A \emph{polytope} $P\subset \R^d$ is the convex hull of a set of finitely many points in $\R^d$. We will always assume that $P$ is $d$-dimensional. A \emph{proper face} of $P$ is the intersection of $P$ with a supporting hyperplane. A polytope $P$ is called \emph{simplicial} if all of its proper faces are geometric simplices, \ie convex hulls of affinely independent points. We identify each face of a simplicial polytope $P$ with the set of its vertices. The \emph{boundary complex} of $P$, denoted $\partial P$, is then the simplicial complex consisting of the empty set along with the vertex sets of proper faces of $P$. A polytope $P$ is called \emph{cs} if $P=-P$; in this case, the complex $\partial P$ is a cs simplicial complex w.r.t.~the natural involution. An important example is $\partial\C^*_d$ --- the boundary complex of a $d$-cross-polytope $\C^*_d\coloneqq \conv(\pm p_1,\pm p_2,\ldots, \pm p_d)$, where $p_1,\ldots,p_d$ are affinely independent points in $\R^d\backslash\{0\}$. As an abstract simplicial complex, $\partial\C^*_d$ is the $d$-fold suspension of $\{\emptyset\}$. It is easy to check that $h_j(\partial \C_d^*)=\binom{d}{j}$ for all $0\leq j\leq d$, and so $g_j(\partial \C_d^*)=\binom{d}{j}-\binom{d}{j-1}$ for all $1\leq j\leq d/2$.


The free involution $\alpha$ on a cs complex $\Delta$ induces the free involution on $X$ via $\alpha(x_v)=x_{-v}$, which in turn induces a $\Z/2\Z$-action on $\R[X]$ and $\R[\Delta]$. For any $\R$-vector space $W$ endowed with such an action $\alpha$, one has $W=W^+\oplus W^-$, where $W^+\coloneqq \{w\in W: w=\alpha(w)\}$ and $W^-\coloneqq \{w\in W: w=-\alpha(w)\}$. Thus,
$\R[\Delta]_i=\R[\Delta]^+_i\oplus \R[\Delta]^-_i$. As $\R[\Delta]^+_{i}\cdot\R[\Delta]^-_j \subseteq \R[\Delta]^{-}_{i+j}$, and similar inclusions hold for all choices of plus and minus signs, it follows that $\R[\Delta]$ has an $(\N\times \Z/2\Z)$-grading. 

Let $\Delta$ be a cs simplicial complex with an involution $\alpha$, and let $\Theta=\theta_1,\ldots,\theta_\ell$ consist of linear forms that are homogeneous w.r.t.~the $(\N\times \Z/2\Z)$-grading. Since $\alpha(I_\Delta+(\Theta))=I_\Delta+(\Theta)$ and since for any $w,w'\in \R[X]_i$, $\langle\alpha(w),\alpha(w')\rangle=\langle w,w'\rangle$, where $\langle-,-\rangle$ is the inner product from~\cite[Section 3]{Lee96} used to define the isomorphism $\Phi_i$ between $\R(\Delta, \Theta)_i$ and $\Stress(\Delta, \Theta)_i$, it follows that $\alpha$ also acts on $\Stress(\Delta, \Theta)_i$ and that this action commutes with $\Phi_i$. Hence, $\Stress(\Delta, \Theta)_i=\Stress(\Delta, \Theta)_i^+ \oplus \Stress(\Delta, \Theta)_i^-$, where the subspaces $\Stress(\Delta, \Theta)^+_i$ and $\Stress(\Delta, \Theta)^-_i$ of $\Stress(\Delta,\Theta)_i$ are isomorphic (as vector spaces) to $\R(\Delta, \Theta)^+_i$ and $\R(\Delta, \Theta)^-_i$, resp. We refer to the elements of $\Stress(\Delta, \Theta)^+_i$ as \emph{symmetric $i$-stresses}. 

